\section{Empirical Methodology}
\label{sec:methodology}

The current evaluation is designed as a transparent empirical record rather than
a controlled ablation study. Recent research-agent benchmarks instead hold tasks,
evaluators, and time or resource budgets fixed to support direct comparisons
~\citep{chan2024mlebench,wijk2024rebench,wang2026aarri,lupidi2026airsbench}. The
The analysis asks whether the deployed runtime has produced competitive
artifacts across heterogeneous settings and how broadly it has been used.

\subsection{Research Questions}

We organize the evaluation around two questions.

\paragraph{RQ1: Cross-domain performance.}
What task-native results has \system{} produced across distinct benchmark arenas,
and how do those results compare with the human or paper-reported references
published for each arena?

\paragraph{RQ2: Research breadth.}
How many de-duplicated manuscripts and drafts has the research pipeline produced,
and how are they distributed across research programs?

The questions deliberately avoid causal language. Although all six runs use
GPT-5.5 through the Codex backend, the records do not share one task interface,
random-seed policy, evaluator, or hardware protocol. Consequently, they can
establish operational breadth and task-level outcomes, but not the isolated
contribution of role separation, checkpointing, review, or any other individual
mechanism.

\subsection{Benchmark Definitions}

Following the task-definition style used by Arbor~\citep{jin2026arbor}, we define
each arena by its research artifact, evaluator, native metric, and metric direction.
The common research-agent configuration is GPT-5.5 through Codex; the target
artifact and evaluator differ by arena.

\paragraph{NVIDIA SOL-ExecBench.}
SOL-ExecBench evaluates correctness-preserving optimization of real GPU kernels
against workload-specific hardware speed-of-light models~\citep{lin2026solexecbench}.
Candidate implementations must first match the reference output; performance is
then assessed relative to the hardware limit. The reported Argus result covers 101
B200 kernels and uses global rank, first-place finishes, and top-three placements as
the public summary; better rank and more high placements are better.

\paragraph{nanochat on B200.}
This arena instantiates the nanochat autoresearch task on one B200 GPU
~\citep{karpathy_nanochat,karpathy_autoresearch}. The editable artifact is the
training program, and each candidate receives a fixed five-minute training window.
The evaluator reports validation bits per byte (BPB), a tokenizer-independent
compression metric; lower BPB indicates a better trained model. The public record
summarizes 426 autonomous attempts.

\paragraph{nanochat on H100.}
The H100 arena uses the same five-minute training objective and lower-is-better BPB
metric, but executes on one H100 and therefore constitutes a separate
hardware-specific optimization problem. The public result summarizes a search over
37 implemented mechanisms. B200 and H100 values are compared only with references
measured under their corresponding hardware protocol, not directly with each other.

\paragraph{nanoGPT speedrun.}
NanoGPT-Bench measures the wall-clock time required for an eight-H100 training run
to reach the fixed validation-loss threshold of $3.28$~\citep{jordan_nanogptbench}.
A candidate is valid only when it reaches that target; among valid candidates,
lower training time is better. The reported value is the mean of $N{=}10$ verifier
runs for the selected artifact.

\paragraph{AARRI-Bench.}
AARRI-Bench evaluates models and agentic harnesses on granular tasks from the
research lifecycle, including research judgment and professional research
practice~\citep{wang2026aarri}. The Argus evaluation uses 82 research-intern tasks
and reports the number and percentage solved. Higher solve rate is better; the
paper-reported best is used as the external comparison.

\paragraph{Math-Reasoning Data Synthesis (Arbor AO suite).}
This row is the Math-Reasoning Data Synthesis task from Arbor's Autonomous
Optimization suite, not a generic benchmark named ``Arbor''~\citep{jin2026arbor}.
The editable artifact is a pipeline that generates AIME-style reasoning problems.
Arbor evaluates generated problems with a GPT-5.5 ReAct solver and reports the mean
$\mathrm{pass@4}-\mathrm{pass@1}$ gap: a larger gap rewards problems that are not
solved immediately but become solvable with additional attempts. The published
comparison uses Argus 28.0, Arbor 20.83, Claude Code 8.33, and Codex 6.25.

\subsection{Evaluation Sources}

The benchmark record is a snapshot of the six public result cards on the Argus
results page~\citep{argus_results}. It was retrieved on July 14, 2026 with the final
URL, HTTP status, and raw-page SHA-256 stored in
\srcpath{evidence/website\_results.json}. Each card preserves the arena, protocol,
reported result, comparison text, source URL, and corroboration status.

The research-output record is a de-duplicated inventory of the public Argus paper
collection~\citep{argus_research}, stored in
\srcpath{evidence/paper\_inventory.json}. It distinguishes manuscripts from drafts,
removes duplicate versions, and groups artifacts into six programs. The inventory
is not a publication-acceptance record and is not compared with the paper output of
other agents.

The two quantitative figures in Section~\ref{sec:results} are generated
deterministically from these JSON files. Structural diagrams are explanatory
illustrations; they are not empirical plots. Figure provenance is summarized in
Appendix~\ref{app:evidence-map}.

\paragraph{Provenance scope.}
Two rows---nanochat B200 and the nanoGPT speedrun---add artifact identifiers,
verifier excerpts, and content digests to the public result snapshot. The other four
rows preserve the public source and page hash but do not yet provide a complete
artifact bundle. This distinction limits reproducibility and is documented in
Appendix~\ref{app:evidence-map}; it is not an additional performance dimension.

\subsection{Metrics and Comparison Policy}

The six arenas use incompatible metrics: global rank, bits per byte, wall-clock
time, solve rate, and a source-reported gap value. We do not normalize these values
or average them into a single score. Each result is shown in its native unit with
its own protocol and comparison. Lower is better for bits per byte and wall-clock
time. For the Math-Reasoning Data Synthesis comparison, the Arbor task definition
identifies the mean $\mathrm{pass@4}-\mathrm{pass@1}$ gap as higher-is-better
~\citep{jin2026arbor}.

Human records and paper-reported bests are used when the public card provides them.
The SOL-ExecBench row instead uses global rank as its headline and retains the
source's system comparison as secondary context. No significance test is introduced
where the source record does not provide the underlying repeated measurements.

\subsection{Portfolio Counting Policy}

The paper inventory counts a research artifact once after duplicate versions are
removed. ``Manuscript'' and ``draft'' are artifact states, not peer-review outcomes.
The headline total therefore answers RQ2---how much research material the pipeline
produced across programs---but does not measure novelty, correctness, acceptance,
or scientific impact. Those questions require external review beyond the scope of
the current records.
